%Version 3.1 December 2024
% See section 11 of the User Manual for version history
%
%%%%%%%%%%%%%%%%%%%%%%%%%%%%%%%%%%%%%%%%%%%%%%%%%%%%%%%%%%%%%%%%%%%%%%
%%                                                                 %%
%% Please do not use \input{...} to include other tex files.       %%
%% Submit your LaTeX manuscript as one .tex document.              %%
%%                                                                 %%
%% All additional figures and files should be attached             %%
%% separately and not embedded in the \TeX\ document itself.       %%
%%                                                                 %%
%%%%%%%%%%%%%%%%%%%%%%%%%%%%%%%%%%%%%%%%%%%%%%%%%%%%%%%%%%%%%%%%%%%%%

%%\documentclass[referee,sn-basic]{sn-jnl}% referee option is meant for double line spacing

%%=======================================================%%
%% to print line numbers in the margin use lineno option %%
%%=======================================================%%

%%\documentclass[lineno,pdflatex,sn-basic]{sn-jnl}% Basic Springer Nature Reference Style/Chemistry Reference Style

%%=========================================================================================%%
%% the documentclass is set to pdflatex as default. You can delete it if not appropriate.  %%
%%=========================================================================================%%

%%\documentclass[sn-basic]{sn-jnl}% Basic Springer Nature Reference Style/Chemistry Reference Style

%%Note: the following reference styles support Namedate and Numbered referencing. By default the style follows the most common style. To switch between the options you can add or remove “Numbered” in the optional parenthesis. 
%%The option is available for: sn-basic.bst, sn-chicago.bst%  
 
%%\documentclass[pdflatex,sn-nature]{sn-jnl}% Style for submissions to Nature Portfolio journals
%%\documentclass[pdflatex,sn-basic]{sn-jnl}% Basic Springer Nature Reference Style/Chemistry Reference Style
\documentclass[pdflatex,sn-mathphys-num]{sn-jnl}% Math and Physical Sciences Numbered Reference Style
%%\documentclass[pdflatex,sn-mathphys-ay]{sn-jnl}% Math and Physical Sciences Author Year Reference Style
%%\documentclass[pdflatex,sn-aps]{sn-jnl}% American Physical Society (APS) Reference Style
%%\documentclass[pdflatex,sn-vancouver-num]{sn-jnl}% Vancouver Numbered Reference Style
%%\documentclass[pdflatex,sn-vancouver-ay]{sn-jnl}% Vancouver Author Year Reference Style
%%\documentclass[pdflatex,sn-apa]{sn-jnl}% APA Reference Style
%%\documentclass[pdflatex,sn-chicago]{sn-jnl}% Chicago-based Humanities Reference Style

%%%% Standard Packages
%%<additional latex packages if required can be included here>

\usepackage{graphicx}%
\usepackage{multirow}%
\usepackage{amsmath,amssymb,amsfonts}%
\usepackage{amsthm}%
\usepackage{mathrsfs}%
\usepackage[title]{appendix}%
\usepackage{xcolor}%
\usepackage{textcomp}%
\usepackage{manyfoot}%
\usepackage{booktabs}%
\usepackage{algorithm}%
\usepackage{algorithmicx}%
\usepackage{algpseudocode}%
\usepackage{listings}%
\usepackage{mathrsfs}%
%%%%

%%%%%=============================================================================%%%%
%%%%  Remarks: This template is provided to aid authors with the preparation
%%%%  of original research articles intended for submission to journals published 
%%%%  by Springer Nature. The guidance has been prepared in partnership with 
%%%%  production teams to conform to Springer Nature technical requirements. 
%%%%  Editorial and presentation requirements differ among journal portfolios and 
%%%%  research disciplines. You may find sections in this template are irrelevant 
%%%%  to your work and are empowered to omit any such section if allowed by the 
%%%%  journal you intend to submit to. The submission guidelines and policies 
%%%%  of the journal take precedence. A detailed User Manual is available in the 
%%%%  template package for technical guidance.
%%%%%=============================================================================%%%%

%% as per the requirement new theorem styles can be included as shown below
\theoremstyle{thmstyleone}%
\newtheorem{theorem}{Theorem}%  meant for continuous numbers
%%\newtheorem{theorem}{Theorem}[section]% meant for sectionwise numbers
%% optional argument [theorem] produces theorem numbering sequence instead of independent numbers for Proposition
\newtheorem{proposition}[theorem]{Proposition}% 
%%\newtheorem{proposition}{Proposition}% to get separate numbers for theorem and proposition etc.

\theoremstyle{thmstyletwo}%
\newtheorem{example}{Example}%
\newtheorem{remark}{Remark}%

\theoremstyle{thmstylethree}%
\newtheorem{definition}{Definition}%

\raggedbottom
%%\unnumbered% uncomment this for unnumbered level heads

\begin{document}

\title[Article Title]{PEARL: Parameter-Efficient Adaptation with Retrieval-Augmented Learning for Molecular Property Prediction}

%%=============================================================%%
%% GivenName	-> \fnm{Joergen W.}
%% Particle	-> \spfx{van der} -> surname prefix
%% FamilyName	-> \sur{Ploeg}
%% Suffix	-> \sfx{IV}
%% \author*[1,2]{\fnm{Joergen W.} \spfx{van der} \sur{Ploeg} 
%%  \sfx{IV}}\email{iauthor@gmail.com}
%%=============================================================%%

\author[1]{\fnm{Kavya} \sur{Agarwal}}

\author[1]{\fnm{Sukrit} \sur{Gupta}}

\author[2]{\fnm{Halima} \sur{Bensmail}}

\author*[2]{\fnm{Raghvendra} \sur{Mall}}\email{ramall@hbku.edu.qa}


\affil[1]{\orgdiv{Center for Applied Research in Data Science}, 
\orgname{Indian Institute of Technology Ropar}, 
\orgaddress{\city{Rupnagar}, \postcode{140001}, 
\state{Punjab}, \country{India}}}

\affil*[2]{\orgdiv{Qatar Computing Research Institute}, 
\orgname{Hamad Bin Khalifa University}, 
\orgaddress{\city{Doha}, \country{Qatar}}}

%%================================%%
%% structured abstract %%
%%================================%%

\abstract{
\textbf{Background}: Molecular property prediction requires representations that capture both the sequential chemistry encoded in SMILES strings and the three-dimensional geometry governing binding and toxicity. No systematic benchmark has directly compared parameter-efficient adaptation of 1D chemical language models (CLMs) against a SE(3)-equivariant 3D model across diverse molecular tasks, or examined whether retrieval-augmented context can bridge the gap between these representation classes.

\textbf{Methods:} We introduce PEARL (\textbf{P}arameter-\textbf{E}fficient \textbf{A}daptation with \textbf{R}etrieval-augmented \textbf{L}earning), which (i)~benchmarks Low-Rank Adaptation (LoRA) uniformly across three CLM families - ChemBERTa, MolFormer, and Uni-Mol - on four molecular property datasets, and (ii)~introduces a Retrieval-Augmented Feature Enhancement (RAFE) pipeline that enriches predictions with label-free physicochemical neighborhood context from the ZINC-250k knowledge base.

\textbf{Results:} Uni-Mol's SE(3)-equivariant architecture achieves the best BACE inhibitor classification (MCC$=$0.623), an $8\%$ absolute gain over the best SMILES-based model, confirming that 3D geometry provides an inductive advantage for binding-affinity tasks. On BBBP, RAFE with MolFormer achieves MCC$=$0.798, matching end-to-end LoRA while providing interpretable chemical-neighborhood context. On the five-class Flavor task, RAFE reaches MCC$=$0.804, though, statistically indistinguishable from end-to-end LoRA (MCC$=$0.801). All models achieve $75$-$97\%$ parameter reduction via LoRA, matching or exceeding published baselines including GROVER-Large and Mole-BERT.

\noindent\textbf{Availability:}
Code is available at \url{https://github.com/raghvendra5688/PEARL}.}

\keywords{PEARL, Parameter-Efficient Finetuning, Chemical Language Models, Low-Rank Adaptation, 3D Molecular Representations, Retrieval-Augmented Generation, Molecular Properties, Chemoinformatics}
%%\pacs[JEL Classification]{D8, H51}

%%\pacs[MSC Classification]{35A01, 65L10, 65L12, 65L20, 65L70}

\maketitle
\section{Introduction}\label{sec1}
The study of molecular structure-property relationships has been central to advances in computational drug discovery. Although traditional quantitative structure-activity relationship (QSAR) methods rely on handcrafted molecular descriptors, deep learning techniques, particularly transformer-based architectures~\cite{vaswani2017attention}, enable end-to-end learning of molecular representations directly from their Simplified Molecular Input Line Entry System (SMILES) sequences. This has catalysed the development of chemical language models (CLMs) pretrained on massive molecular corpora, now rivaling experimental approaches in predicting properties such as toxicity, solubility, and bioactivity~\cite{ahmad2022chemberta2,huang2020deeppurpose}.

Despite their promise, two critical challenges persist: (1) CLMs process SMILES as flat token sequences and therefore lack explicit inductive biases for three-dimensional molecular geometry, which is fundamental to binding affinity and enzyme inhibition; and (2) the computational cost of fine-tuning $\mathcal{O}(10^8)$-parameter CLMs for each specialized task remains prohibitive in resource-constrained settings.

This work addresses both challenges through a systematic study of three model families: MolFormer~\cite{ross2022molformer}, ChemBERTa-2~\cite{ahmad2022chemberta2}, and Uni-Mol~\cite{zhou2023unimol}, a SE(3)-equivariant Transformer that operates directly on explicit 3D atomic coordinates. We apply Low-Rank Adaptation (LoRA)~\cite{hu2021lora} to all three families, dramatically reducing trainable parameters while preserving or exceeding the performance of full finetuning.

In addition, we introduce a \emph{Retrieval-Augmented Feature Enhancement} (RAFE) framework that uses task-specific finetuned embeddings to retrieve nearest chemical neighbors from ZINC-250k~\cite{gomez2018automatic}, a 249,455-molecule drug-like library, via GPU-accelerated FAISS indices~\cite{johnson2021faiss}. The physico-chemical neighborhood context derived from these retrievals (logP, QED, and SAS distributions) provides an orthogonal, label-free signal that complements learned model representations.

Our key contributions are:
\begin{itemize}
  \item A hybrid framework, \textbf{PEARL}, that supports three CLM families (ChemBERTa, MolFormer, Uni-Mol) with LoRA finetuning, achieving $75$-$96\%$ parameter reduction across all model sizes.
  \item A \textbf{3D geometry-aware} molecular property predictor via Uni-Mol LoRA, which yields the best BACE inhibitor classification performance (MCC=$0.623$) across all methods studied.
  \item A \textbf{Retrieval-Augmented Feature Enhancement} (RAFE) framework using ZINC-250k as an external, leakage-free knowledge base; PEARL achieves MCC=$0.804$ on five-class Flavor prediction versus end-to-end (E2E) LoRA's MCC=$0.801$ and MCC=$0.798$ on BBBP, matching E2E LoRA while providing chemical-neighborhood context.
  \item Comprehensive ablation across four molecular property prediction datasets, two loss functions, and three tree-based ensemble classifiers, highlighting task-dependent optimal strategies.
  \item We made all the code publicly available for the community at \url{https://github.com/raghvendra5688/PEARL}.
\end{itemize}

Figure~\ref{fig:workflow} illustrates the complete PEARL workflow. The molecular dataset is processed by one of three CLM backbones. LoRA adapters augment only the attention projection layers, keeping all other pretrained weights frozen. For Uni-Mol, 3D conformers (SE(3)) are pre-generated via ETKDGv3 + MMFF94 and cached.
The finetuned model provides both end-to-end predictions and task-specific embeddings. These embeddings are optionally concatenated with engineered PC features and with ZINC-250k RAFE features retrieved via FAISS indices, and the resulting feature vector is fed to Optuna-tuned tree-based classifiers.

\begin{figure*}[!ht]
  \centering
  \includegraphics[width=0.9\textwidth]{Figures/PEARL_Workflow_v2.pdf}
  \caption{PEARL workflow. Molecules are processed by one of three CLM families (ChemBERTa, MolFormer, Uni-Mol) with LoRA adapters. Task-specific embeddings are fused with molecular PC features and ZINC-250k retrieval features, then passed to tree-based classifiers. Uni-Mol additionally encodes 3D ETKDGv3 conformers via SE(3) representation.}
  \label{fig:workflow}
\end{figure*}

\section{Related Work}\label{related_work}

\subsection{Chemical language models.} \label{CLMs}
ChemBERTa~\cite{chithrananda2020chemberta} and ChemBERTa-2~\cite{ahmad2022chemberta2} adapted BERT-style masked language modelling to SMILES, establishing the template for CLM pretraining on large molecular corpora.
MolFormer~\cite{ross2022molformer} scaled this to $1.1\times10^9$ molecules using linear attention.
Uni-Mol~\cite{zhou2023unimol} departed from SMILES by operating on explicit 3D atomic coordinates via an SE(3)-equivariant transformer, achieving state-of-the-art results on MoleculeNet through full finetuning.
Multi-modal models such as MolT5~\cite{edwards2022molt5}, nach0~\cite{livne2024nach0}, and Omni-Mol~\cite{hu2025omnimol} further integrate SMILES, graphs, and natural language, but systematic evaluation under parameter-efficient training conditions remains limited.

\subsection{Parameter-efficient finetuning (PEFT) for molecular CLMs.} \label{PEFT}
LoRA~\cite{hu2021lora} introduced low-rank weight updates for large language models; PEFT methods have since been evaluated for molecular tasks.
Balne et al.~\cite{balne2024peft} and Xu et al.~\cite{xu2023peftreview} survey PEFT strategies broadly, including adapter layers, prefix tuning, and LoRA, while
Fu et al.~\cite{fu2023effectiveness} assess their effectiveness in NLP.
Most closely related to PEARL is EffiChem~\cite{harod2025effichem}, a prior work that demonstrates efficient LoRA-based adaptation of 1D CLMs (ChemBERTa, MolFormer) for molecular property prediction.
PEARL expands EffiChem in two directions: (i)~it adds Uni-Mol, a 3D SE(3)-equivariant model, enabling the first direct LoRA comparison between 1D SMILES-based and 3D geometry-aware CLMs on identical scaffold splits; and (ii)~it introduces the RAFE pipeline, which augments frozen CLM embeddings with label-free chemical neighborhood features from an external knowledge base.
To our knowledge, no prior work combines PEFT with leakage-free retrieval augmentation across both 1D and 3D molecular representations.

\subsection{Retrieval-augmented methods in molecular ML.} \label{RAG}
Generative RAG~\cite{lewis2020rag} conditions a language model on retrieved documents; molecular analogues typically retrieve labeled training-set molecules as auxiliary signals, introducing leakage and a small retrieval pool.
KPGT~\cite{liu2022kpgt} injects knowledge-graph attributes into graph-transformer pretraining but retrieval occurs at pretraining time from the training corpus itself.
PEARL's RAFE differs: it retrieves from a fixed external corpus (ZINC-250k) at inference time, extracts label-free physicochemical statistics as tree-classifier features, and applies to any CLM family without retraining.

\subsection{GNNs and 3D representations.} \label{Graph}
Chemprop~\cite{yang2019chemprop}, AttentiveFP~\cite{xiong2020attentivefp}, GROVER~\cite{rong2020grover}, and Mole-BERT~\cite{xia2023molecbert} demonstrate the strength of graph-based molecular representations.
PEARL includes these as external baselines in Supplementary Section~S5 and treats Uni-Mol as the 3D representative, benchmarking all approaches under identical scaffold splits to isolate the contribution of geometry-aware representations.
%%%%%%%%%%%%%%%%%%%%%%%%%%%%%%%%%%%%%%%%%%%%%%%%%%%%%%%%%%%%%%%%%%

\section{Materials and Methods} \label{Materials_Methods}

\subsection{Data Collection and Curation} \label{data_collection}
The \textbf{Blood-Brain Barrier Penetration (BBBP)} dataset was originally curated in~\cite{martins2012bayesian} for predicting barrier permeability and contains over $2,000$ compounds with binary labels, integrated into the MoleculeNet~\cite{wu2018moleculenet} benchmark. \textbf{ClinTox} is part of the MoleculeNet benchmark~\cite{wu2018moleculenet}. The dataset contains 1,476 drugs that either failed clinical trials due to toxicity or succeeded, along with FDA approval status. FDA-approved drugs originate from the SWEETLEAD database~\cite{novick2013sweetlead}, while toxicity-failed candidates derive from the Aggregate Analysis of Clinical Trials (AACT).

The \textbf{Flavor Analysis and Recognition Transformer (FART)} dataset, introduced in~\cite{zimmermann2024fart}, is a multi-class taste prediction corpus of 15,025 compounds annotated as \textit{bitter}, \textit{sour}, \textit{sweet}, \textit{umami}, and \textit{undefined}. The \textbf{BACE} dataset is part of MoleculeNet~\cite{wu2018moleculenet} and contains experimentally measured binding results for 1,513 inhibitors of human $\beta$-secretase~1 (BACE-1), with binary activity labels. Dataset splits were curated using the DeepChem reference implementation~\cite{ramsundar2019deep}. BBBP and ClinTox datasets are sourced from the IBM MolFormer repository and underwent validity checks, duplicate removal, and class-ambiguity filtering.

\subsection{Data Splits and Class Distribution} \label{data_splits_distribution}
BBBP, ClinTox, and BACE employ scaffold-based splits~\cite{ramsundar2019deep} using an 80:10:10 train/validation/test ratio. BBBP exhibits moderate class imbalance (78.8\% permeable in training). ClinTox faces severe imbalance: non-toxic compounds dominate at $92.3\%$. The Flavor dataset uses the curated splits provided in~\cite{zimmermann2024fart}, with pronounced multi-class skew (sweet $41\%$, bitter $32\%$, sour $11\%$, umami $4\%$, undefined $12\%$). BACE is the most balanced, with $54.3\%$ inactive and $45.7\%$ active compounds.

\subsection{Pretrained Chemical Language Models} \label{Pre_CLMs}
\textbf{ChemBERTa-2}~\cite{ahmad2022chemberta2} builds upon ChemBERTa~\cite{chithrananda2020chemberta} via a RoBERTa-style architecture trained on $77$ million canonical SMILES from PubChem~\cite{kim2023pubchem}. Two pretraining strategies are used: Masked Language Modeling (MLM) and Multi-Task Regression (MTR) on $200$ molecular
property tasks. We use the two 77M-parameter variants, denoted \textbf{ChemBERTa-77M~MLM} and \textbf{ChemBERTa-77M~MTR}, each producing 384-dimensional embeddings.

\textbf{MolFormer}~\cite{ross2022molformer} is a scalable CLM that learns representations from SMILES without 3D geometry. It uses a transformer with linear attention (FAVOR+)~\cite{choromanski2021rethinking}, rotary positional embeddings~\cite{heo2024rotary}, and efficient distributed training~\cite{isard2009distributed}, enabling pretraining on $1.1\times 10^{9}$  molecules from PubChem~\cite{kim2023pubchem} and ZINC20~\cite{irwin2020zinc20}. We use the \textbf{MolFormer-XL-10pct-both} variant ($\approx 46$M parameters, pretrained on $10\%$ of both datasets) with 768-dimensional embeddings.

\textbf{Uni-Mol}~\cite{zhou2023unimol} is a universal 3D molecular pre-trained model based on a SE(3)-equivariant Transformer with 15 attention layers. Unlike SMILES-based CLMs, Uni-Mol accepts \emph{explicit 3D atomic coordinates} generated via RDKit ETKDGv3 conformer embedding followed by MMFF94 force-field minimization.
This geometry-aware design gives Uni-Mol inductive biases suited to structure-dependent properties such as receptor binding affinity. For PEARL, each molecule is represented as a \textbf{2560-dimensional hybrid vector}: a 512-dimensional CLS token from the SE(3)-Transformer concatenated with 2048-dimensional Morgan ECFP4 fingerprints (radius $= 2$), combining 3D structural awareness with explicit 2D topological information.

\subsection{Low-Rank Adaptation (LoRA)}\label{sec:lora}
Low-Rank Adaptation (LoRA)~\cite{hu2021lora} introduces two trainable low-rank matrices $\mathbf{A} \in \mathbb{R}^{d \times r}$ and $\mathbf{B} \in \mathbb{R}^{r \times d}$ alongside frozen pretrained weight matrices $\mathbf{W} \in \mathbb{R}^{d \times d}$ in the attention layers, where $r \ll d$ is the rank hyperparameter. The weight update is: 
\begin{equation}
  \mathbf{W}' = \mathbf{W} + \Delta\mathbf{W}, \quad \Delta\mathbf{W} = \mathbf{A} \cdot \mathbf{B}.
\end{equation}
Matrix $\mathbf{A}$ is initialized from $\mathcal{N}(0, \sigma^2)$ and $\mathbf{B}$ from zeros, ensuring that adaptation begins from the pretrained representation.

\subsubsection{LoRA for ChemBERTa and MolFormer}
LoRA adapters are applied to the query ($\mathbf{W}^Q$), key ($\mathbf{W}^K$), and value ($\mathbf{W}^V$) projection matrices of each attention layer. All other parameters are frozen. Hyperparameter search covers rank $r \in \{4, 8, 16, 32, 64\}$, LoRA scaling $\alpha \in \{8, 16, 32, 64, 128\}$, and dropout $\in \{0.0, 0.1, 0.2\}$.

\subsubsection{LoRA for Uni-Mol} The version of Uni-Mol employs a fused attention projection, where the query, key, and value transformations are combined into a single linear layer (\texttt{in\_proj}) rather than separate $\mathbf{W}^Q$, $\mathbf{W}^K$, and $\mathbf{W}^V$. Consequently, LoRA is applied directly to this unified projection module across all self-attention layers.
Specifically, LoRA adapters are injected into the \texttt{in\_proj} weights (and \texttt{out\_proj} when available), enabling parameter-efficient fine-tuning without modifying the full backbone. The rank $r$ of the low-rank decomposition is selected from $\{4, 8, 16, 32\}$, while the scaling factor $\alpha$ is chosen from $\{8, 16, 32, 64\}$. 

\begin{figure}[!htb]
  \centering
  \includegraphics[width=0.8\textwidth]{Figures/LoRA_FineTuning.pdf}
  \caption{Low-Rank Adaptation (LoRA) Architecture. A low-rank adaptation framework fine-tunes frozen transformer weights (Q/K/V layers) using trainable low-rank matrices (A, B), enabling efficient adaptation to task-specific molecular datasets (ClinTox, BBBP, BACE, Flavor).} %with lightweight classification heads for binary and multiclass property prediction.}
  \label{fig:lora-architecture}
\end{figure}
Conformers for all dataset molecules are generated once before the hyperparameter sweep and cached, avoiding redundant ETKDGv3 calls during hyperparameter optimization. We use necessary flags to prevent conformer generation failures for metal-coordinates SMILES. Figure \ref{fig:lora-architecture} illustrates how LoRA finetuning is implemented within our framework.

\subsubsection{Training details}
All models use AdamW with cosine-annealing learning rate schedules. The encoder (LoRA adapters) receives the base learning rate; classification head receives $5\times$ the base rate. Gradient clipping ($\text{max\_norm} = 1.0$) is applied over all parameters. Early stopping monitors validation MCC with patience of 5 epochs. Bayesian hyperparameter optimization with WandB sweeps ($20$-$30$ trials per dataset) optimizes for validation MCC.

\subsection{Loss Functions for Class Imbalance} \label{loss_functions}
\noindent Two loss functions address the severe class imbalance prevalent in molecular datasets, implemented via a custom extension of the HuggingFace \texttt{Trainer}.

\noindent \textbf{Weighted cross-entropy (WL):}\\
$\mathcal{L}_{\text{WL}} = -\sum_{i=0}^{C-1} w_i y_i \log\hat{y}_i$, where $w_i = 1 - f_i$ (inverse class frequency). This gives the minority class higher loss contribution and is most effective for datasets like BACE and Flavor.

\noindent\textbf{Focal loss (FL):}\\
$\mathcal{L}_{\text{FL}} = -\alpha (1 - \hat{p}_i)^\gamma \log\hat{p}_i$ with $\gamma = 2$, which down-weights well-classified examples to focus learning on hard-to-classify boundary cases. This excels on heavily imbalanced and challenging tasks such as BBBP and ClinTox.

\subsection{PEARL Architecture} \label{PEARL}
PEARL defines four operational modes that form a progressive evaluation ladder:

\textbf{(a)~End-to-End (E2E) LoRA finetuning:}
The CLM backbone (ChemBERTa, MolFormer, or Uni-Mol) with LoRA adapters and a task-specific classification head is optimized end-to-end for each downstream task. For binary tasks, the head is a dense layer with sigmoid activation;
for multi-class tasks, a dense layer with softmax.

\textbf{(b)~Finetuned Embedding Prediction (FT~Embed):}
Task-specific embeddings are extracted from the finetuned model. For MolFormer, the CLS token passes through two dense layers with GELU activation (768-dim output). For ChemBERTa, the CLS token passes through a dense layer with tanh activation (384-dim). For Uni-Mol, the CLS token is concatenated with ECFP4 fingerprints (2560-dim).
These embeddings train Optuna-tuned \cite{akiba2019optuna} tree-based classifiers (XGBoost / LightGBM / CatBoost).

\textbf{(c)~FT~Embed $+$ PC (Physicochemical Feature Fusion).}
The FT~Embed vector is concatenated with 473~domain-specific physicochemical (PC) features: 148 RDKit molecular descriptors, 30 graph-based features (NetworkX v3.5), 128-bit Morgan fingerprints, and 167-bit MACCS keys (RDKit v2025.09.1). This combined vector trains the same tree-based classifiers without any retrieval from external databases. A full description of engineered features is provided in Supplementary Table~S1.

\textbf{(d)~Retrieval-Augmented Feature Enhancement (RAFE).}
RAFE augments the FT~Embed $+$ PC feature vector with ZINC-250k chemical neighborhood features (69-dim; see Section~\ref{sec:rafe}), yielding the full feature vector [$\text{CLM embeddings} \mid \text{PC features (473-dim)} \mid \text{ZINC-250k RAFE features} (69\text{-dim})$], passed to the similar Optuna-tuned tree-based classifiers.

\subsection{RAFE} \label{sec:rafe}
\subsubsection{Design Rationale and External Knowledge Base} \label{sec:zinc_rationale}
Naive retrieval-augmented approaches retrieve nearest neighbors from the task-specific training set and use their labels as auxiliary features. This introduces three problems: a small retrieval pool ($\sim$1,200-12,000 molecules), training-set leakage, and circular dependencies between the retrieval space and the prediction target.

\noindent\textbf{Note on terminology:} Unlike generative RAG~\cite{lewis2020rag}, where retrieved passages are injected into a language model prompt, RAFE does not pass retrieved content to any neural model. Retrieved ZINC-250k neighbors serve solely as a source of label-free physicochemical statistics (logP, QED, SAS distributions, structural similarity) that are featurised and concatenated with the molecular embedding before a tree-based classifier. The ``retrieval-augmented'' framing refers to the use of an external knowledge base at inference time, not to in-context generation. 

PEARL addresses all three by using \textbf{ZINC-250k}~\cite{gomez2018automatic} as the retrieval corpus i.e. a curated subset of $249,455$ drug-like molecules from the ZINC database~\cite{irwin2020zinc20} with no overlap with
any training, validation, or test split. Each ZINC-250k molecule carries three physicochemical properties: lipophilicity (logP), Quantitative Estimate of Drug-likeness (QED $\in [0,1]$), and Synthetic Accessibility Score (SAS $\in [1,10]$).

ZINC-250k is chosen as the retrieval corpus for four complementary reasons: 

1. \textbf{Scale and diversity:} $249,455$ drug-like molecules span a broad region of chemical space, providing a dense and representative neighborhood for any query molecule from BACE, BBBP, ClinTox, or Flavor, all datasets have an at most 15,025 compounds.

2. \textbf{Zero leakage:} ZINC-250k has no label overlap with any training, validation, or test split used in this study. Retrieved neighbors cannot carry prediction targets, eliminating the circular dependency that plagues training-set retrieval.

3. \textbf{Label-free property signal:} logP, QED, and SAS are computed directly from SMILES without biological assay data. Their distributions in the local neighborhood of a query molecule characterize its physico-chemical environment in a way that is complementary to the molecule's own PC features and to the CLM's learned representation.

4. \textbf{Model-family independence:} ZINC-250k is embedded independently by each finetuned CLM variant (ChemBERTa, MolFormer, Uni-Mol), yielding model-specific FAISS indices. RAFE features are therefore always expressed in the same geometric space as the query molecule's embedding, ensuring neighborhood statistics remain geometrically coherent with the representations passed to the downstream classifier.

The core insight is that the logP/QED/SAS \emph{distributions of retrieved nearest ZINC neighbors} characterize the local chemical space around a query molecule in a label-free, leakage-free manner that complements the molecule's own PC features:
(i)~BBB-permeable molecules cluster in high-logP, high-QED ZINC regions~(BBBP);
(ii)~BACE inhibitors reside in drug-like, easily-synthesizable ZINC neighborhoods~(BACE);
(iii)~clinically toxic compounds often exhibit unusual logP profiles and heterogeneous physicochemical environments~(ClinTox);
(iv)~bitter compounds cluster in hydrophobic ZINC regions, while sweet/sour compounds occupy hydrophilic regions.

\subsubsection{Vector Indexing with FAISS} \label{Vector_FAISS}
ZINC-250k molecules are embedded using task-specific finetuned model variants across all three CLM families: six SMILES-based variants (ChemBERTa-MLM $\times$ 2 loss, ChemBERTa-MTR $\times$ 2 loss, MolFormer $\times$ 2 loss) and two Uni-Mol variants (FL and WL), yielding 32 embedding matrices in total (8 models $\times$ 4 datasets).

SMILES-based embeddings are 768-dimensional CLS tokens; Uni-Mol embeddings are 2560-dimensional hybrid vectors (512-dim SE(3)-Transformer CLS concatenated with 2048-dim ECFP4 fingerprints). Embeddings are L2-normalized and indexed with FAISS \texttt{GpuIndexFlatIP}~\cite{johnson2021faiss}, a GPU-accelerated exact cosine-similarity
search that achieves sub-2~ms query latency at 249k$\times$768 on a single GPU.
Indices are serialized as CPU \texttt{IndexFlatIP} files for portable, fast retrieval.

\subsubsection{Chemical neighborhood Feature Extraction} \label{feature_extraction}
For every molecule in every split (train, validation, test), each FAISS index is queried with $k \in \{10, 20\}$ nearest ZINC-250k neighbors. The following feature groups are extracted:

\textbf{(a) Physico-chemical neighborhood profile ($\sim$27 features).}
Mean and standard deviation of logP, QED, and SAS across $k = 10$ and $k = 20$ neighbors, plus similarity-weighted means (e.g.\ $\overline{\text{logP}}_w = \sum_i s_i \cdot
\text{logP}_i / \sum_i s_i$ where $s_i$ is cosine similarity to the $i$-th neighbor).

\textbf{(b) Neighborhood density and isolation (4 features).}
Cosine similarity to the nearest ZINC neighbor ($\texttt{zinc\_nearest\_sim}$),
mean and standard deviation of similarities across $k = 10$ neighbors, and Shannon entropy of the quantized similarity distribution. Molecules with low \texttt{zinc\_nearest\_sim} are isolated in chemical space and their predictions are inherently less certain.

\textbf{(c) Embedding centroid features ($\sim$35 features).}
The mean of $k = 10$ neighbor embeddings forms a centroid; PCA is fitted exclusively on training-split centroids and then applied without refitting to validation and test centroids, compressing each to 32 dimensions. Additional features capture cosine similarity of the query to its centroid and the spread (mean and standard deviation of L2 distance of neighbors from the centroid).

\textbf{(d) Tanimoto structural similarity (3 features).}
Tanimoto coefficients of Morgan fingerprints (radius~$= 2$, 128-bit) between the query and its top-1 and top-5 ZINC neighbors, plus the fraction of top-10 neighbors sharing the same Murcko scaffold.

The four groups yield a 69-dimensional RAFE feature vector per molecule. The full augmented feature vector is:
[$\text{CLM embeddings} \mid \text{PC features (473-dim)} \mid \text{ZINC RAFE features}
(69\text{-dim})$].

\begin{figure}[!ht]
  \centering
  \includegraphics[width=0.9\textwidth]{Figures/RAFE_Framework.pdf}
  \caption{Retrieval-Augmented Feature Enhancement (RAFE) pipeline. A query molecule is embedded by a task-specific finetuned CLM; the embedding is used to retrieve $k$ nearest chemical neighbors from the ZINC-250k FAISS index. Four feature groups are extracted: physico-chemical neighborhood profile (logP, QED, SAS distributions; $\sim$27-dim), neighborhood density and isolation (4-dim), embedding centroid features ($\sim$35-dim), and Tanimoto structural similarity (3-dim), yielding a 69-dim RAFE vector concatenated with the CLM embedding and engineered PC features for tree-based classification.}
  \label{fig:rafe}
\end{figure}

\subsection{Tree-based Ensemble Classifiers} \label{Tree_Models}
The final prediction stage employs XGBoost~\cite{chen2016xgboost}, LightGBM~\cite{ke2017lightgbm}, and CatBoost~\cite{prokhorenkova2018catboost}. These algorithms handle high-dimensional mixed-type feature vectors, provide built-in feature importance, and are robust to outliers. Bayesian hyperparameter optimization was performed with Optuna~\cite{akiba2019optuna} ($20$ trials for each configuration) maximizing validation MCC via stratified 5-fold cross-validation. To address class imbalance, class-weighting strategies are incorporated within each model to rebalance the contribution of minority and majority classes during training.

\subsection{Evaluation Metrics} \label{metrics}
Performance is evaluated using accuracy (Acc), area under the ROC curve (AUC), Matthews correlation coefficient (MCC), precision (Prec), recall (Rec), and F1-macro (F1$_\text{mac}$) / F1-micro (F1$_\text{mic}$). MCC is the primary optimization target due to its robustness to class imbalance. %Area under the precision-recall curve (AUPR) is additionally reported.
Given that ClinTox ($n_\text{test}$ = $143$ with 7 positives) and BACE ($n_\text{test} \approx 152$) have small test sets, \textbf{performance differences are interpreted alongside bootstrapped 95\% confidence intervals (10{,}000 stratified resamples, normal-approximation}); Key CIs are reported inline for best-performing configurations; non-overlapping CIs between two models are taken as evidence of a practically statistically significant difference.
%%%%%%%%%%%%%%%%%%%%%%%%%%%%%%%%%%%%%%%%%%%%%%%%%%%%%%%%%%%%

\section{Experimental Results} \label{experimental_results}
Results are organised around the four PEARL operational modes described in Section~2.4. A direct comparison against published external baselines (Chemprop~\cite{yang2019chemprop}, GROVER~\cite{rong2020grover}, AttentiveFP~\cite{xiong2020attentivefp}, Mole-BERT~\cite{xia2023molecbert}, and the original Uni-Mol and MolFormer full-finetuning results) is summarised in Section~\ref{sec:external_comparison}; full method descriptions and ablations are in Supplementary Section~S5. \textbf{Bootstrapped 95\,\% confidence intervals (CI) (10{,}000 stratified resamples) are reported for key metrics of top performing models}.
(i)~\textbf{E2E LoRA}- End-to-end LoRA finetuning of the CLM backbone; 
(ii)~\textbf{FT Embed}- Optuna-tuned tree classifiers on frozen finetuned embeddings;
(iii)~\textbf{FT Embed $+$ PC}- FT Embed augmented with 473-dim physicochemical features, without ZINC-250k retrieval;
(iv)~\textbf{RAFE}- FT Embed $+$ PC further augmented with ZINC-250k neighborhood features. For each mode we report the best-performing CLM variant and, where applicable, the corresponding tree model. Full ablation tables appear in the Supplementary Material.

\subsection{Sequence of Comparisons}\label{sec:comparison_sequence}
The four PEARL modes form an evaluation ladder that isolates the contribution of each component:

\textbf{(i)~E2E LoRA establishes the performance ceiling} achievable by fully differentiable gradient flow through the CLM backbone.  Because LoRA adapters and the classification head are optimized jointly, the model can learn task-specific representations end-to-end. E2E LoRA results are therefore treated as the upper bound for each CLM family and dataset.

\textbf{(ii)~FT Embed quantifies how much embedding quality alone contributes.} Here, the finetuned model is frozen after training; its CLS token embedding is extracted and passed to an Optuna-tuned tree classifier (XGBoost / LightGBM / CatBoost). Any gap between FT Embed and E2E LoRA reflects information that can only be recovered through
end-to-end gradient flow i.e., representations that emerge from joint task-head optimization but are not fully preserved in the frozen CLS token.

\textbf{(iii)~FT Embed $+$ PC isolates the contribution of physicochemical feature engineering.} The frozen embedding is augmented with 473 engineered PC features and passed to the same tree classifier. A gain over FT Embed reflects the additional discriminative power of handcrafted molecular descriptors and fingerprints, independent of any external retrieval.

\textbf{(iv)~RAFE measures the additional value of ZINC-250k chemical neighborhood context.} RAFE augments the FT Embed $+$ PC feature vector with ZINC-250k retrieval features (69-dim), all passed to the same tree classifier family. A gain from FT Embed $+$ PC to RAFE indicates that neighborhood signals from ZINC-250k provide information orthogonal to both the learned CLM representation and the molecule's own physicochemical profile. A lack of gain (or degradation) indicates that the task is primarily driven by geometry-critical or training-distribution-specific signals that retrieval cannot replicate.

This ladder allows us to attribute performance differences to four distinct sources: 
(a)~the quality of the underlying pre-trained CLM, (b)~the fidelity of the frozen embedding, (c)~the discriminative power of engineered PC features, and (d)~the neighborhood context contributed by ZINC-250k. The per-dataset results in the following subsections are discussed in this light.

\subsection{BBBP Dataset Results}
Table~\ref{tab:bbbp} presents the best configuration per modality and CLM family for BBBP. E2E LoRA finetuning of MolFormer-XL with weighted loss achieves optimal overall performance
(MCC$=0.797$, AUC$=0.948$, F1$_\text{mac}=0.898$), matching the best previously reported result for this dataset.

Uni-Mol with weighted loss achieves MCC$=$0.774 and AUC$=$0.943, demonstrating that 3D geometry-aware representations are highly competitive on BBBP, despite this task being conventionally attributed to 2D physicochemical properties such as logP and molecular weight. Uni-Mol FT Embed (WL + CatBoost) achieves MCC$=0.720$, AUC$=0.918$ and is much worse than the more informative MolFormer FT Embed (FL + XGBoost) which achieves MCC$=$0.787 and AUC$=$0.935. Applying RAFE to Uni-Mol embeddings (FL + XGBoost) improves this to MCC$=0.773$, AUC$=$0.928 which is marginally below Uni-Mol E2E LoRA (MCC$=0.774$) while using only tree-based classifiers, and narrowing the gap with MolFormer RAFE (MCC$=0.798$).

The RAFE framework (MolFormer-XL FL + XGBoost) achieves MCC$= 0.798$ (95\,\% CI $[0.743, 0.853]$, SE$= $0.028), \emph{matching} the E2E LoRA result while using a tree-based classifier rather than end-to-end neural finetuning. Importantly, RAFE provides an interpretable neighborhood context signal: the high mean logP of retrieved ZINC-250k neighbors (associated with BBB-permeable molecules) directly explaining the prediction. Adding PC features to finetuned embeddings without RAFE (PC + FT Embed, MolFormer WL + CatBoost: MCC$= 0.773$) is outperformed by both E2E LoRA and RAFE, suggesting that ZINC-250k neighborhood context contributes information orthogonal to the molecule's own PC features.

\begin{table*}[!ht]
\centering
\caption{Best performing configurations for each modality and CLM on the BBBP dataset. $^*$, $^+$ denotes the top and second-best performing method by average metric. RAFE = FT~Embed~+~PC~+~ZINC-250k retrieval features.}
\label{tab:bbbp}
\resizebox{\textwidth}{!}{%
\begin{tabular}{lllllrrrrrrr}
\toprule
Modality & CLM & Loss & ML & AUC & Acc & F1$_\text{mac}$ & F1$_\text{mic}$ & MCC & Prec & Rec \\
\midrule
E2E LoRA & ChemBERTa-77M MLM & FL & --- & 0.923 & 0.880 & 0.851 & 0.865 & 0.739 & 0.889 & 0.851 \\
E2E LoRA & ChemBERTa-77M MTR & FL & --- & 0.921 & 0.870 & 0.843 & 0.854 & 0.715 & 0.873 & 0.843 \\
E2E LoRA$^*$ & MolFormer-XL & WL & --- & 0.948 & \textbf{0.906} & \textbf{0.898} & \textbf{0.906} & 0.797 & 0.901 & \textbf{0.896} \\
%E2E LoRA & Uni-Mol & FL & --- & 0.944 &  0.880 & 0.873 & 0.880 & 0.747 & 0.868 & 0.878 \\
E2E LoRA & Uni-Mol & WL & --- & 0.943 & 0.896 & 0.887 & 0.896 & 0.774 & 0.890 & 0.885 \\
\midrule
FT Embed & ChemBERTa-77M MLM & FL & LightGBM & 0.905 &  0.896 & 0.882 & 0.896 & 0.776 & 0.911 & 0.866 \\
FT Embed & MolFormer-XL & FL & XGBoost & 0.935 & 0.901 & 0.888 & 0.901 & 0.787 & 0.914 & 0.873 \\
%FT Embed & Uni-Mol & FL & XGBoost & 0.917 &  0.865 & 0.842 & 0.865 & 0.709 & 0.889 & 0.823 \\
FT Embed & Uni-Mol & WL & CatBoost & 0.918 & 0.865 & 0.838 & 0.865 & 0.720 & 0.912 & 0.814 \\
\midrule
PC + FT Embed & ChemBERTa-77M MLM & FL & LightGBM & 0.888 & 0.891 & 0.875 & 0.891 & 0.765 & 0.907 & 0.859 \\
PC + FT Embed & MolFormer-XL & FL & LightGBM & 0.934 & 0.901 & 0.887 & 0.901 & 0.789 & 0.920 & 0.870 \\
PC + FT Embed & MolFormer-XL & WL & CatBoost & \textbf{0.951} & 0.896 & 0.884 & 0.896 & 0.773 & 0.902 & 0.872 \\
\midrule
RAFE & ChemBERTa-77M MLM & FL & LightGBM & 0.926 & 0.896 & 0.882 & 0.896 & 0.776 & 0.911 & 0.866 \\
RAFE$^+$ & MolFormer-XL & FL & XGBoost & 0.940 & \textbf{0.906} & 0.895 & \textbf{0.906} & \textbf{0.798} & 0.918 & 0.881 \\
RAFE & MolFormer-XL & WL & CatBoost & 0.946 & 0.885 & 0.872 & 0.885 & 0.750 & 0.889 & 0.861 \\
RAFE & Uni-Mol & FL & XGBoost & 0.928 & 0.891 & 0.872 & 0.891 & 0.773 & \textbf{0.927} & 0.850 \\
\bottomrule
\end{tabular}
}
\end{table*}

\subsection{Flavor Dataset Results}
Table~\ref{tab:flavor} presents results on the multi-class Flavor dataset. E2E LoRA finetuning of MolFormer-XL with weighted loss achieves MCC$=$0.801, F1$_\text{mac}$=$0.809$, AUC$=$0.975, Prec$=$0.802, and Rec$=$0.817.

The RAFE framework (MolFormer-XL WL + CatBoost) achieves MCC$=$0.804 ($95\%$ CI $[0.784, 0.824]$, SE$=$0.010), \emph{statistically indistinguishable from} E2E LoRA (MCC$=$0.801 $95\%$ CI $[0.782, 0.820]$; the overlapping CIs indicate no significant difference at $\alpha = 0.05$), while using tree-based classifiers on the augmented feature vector. This suggests that ZINC-250k neighborhood context contributes orthogonal information for distinguishing multi-class taste profiles: logP differences between the
hydrophobic ZINC neighborhood of bitter compounds versus the hydrophilic neighborhood of sweet/sour compounds complement the CLM's learned representations.

Uni-Mol achieves MCC$=$0.795 with weighted loss (Prec$=$0.747, Rec$=$0.848), competitive with but below MolFormer-XL, indicating that 3D geometric information does not provide a strong additional signal for flavour prediction compared to SMILES-based representations. Uni-Mol FT Embed FL (XGBoost) achieves MCC$=$0.765, Acc$=$0.866, F1$_\text{mac}=$0.731, falling below Uni-Mol E2E LoRA FL (MCC$=$0.790) and WL (MCC$=$0.795), demonstrating that end-to-end gradient flow provides a meaningful benefit for 3D taste representation. Uni-Mol RAFE FL (CatBoost) achieves MCC$=$0.764, Acc$=$0.866, F1$_\text{mac}$=0.714, while WL (CatBoost) achieves MCC$=$0.761, F1$_\text{mac}$=0.737 suggesting WL's minority-class upweighting yields superior macro-averaged balance, particularly for the rarest umami class (AUPR: WL$=$0.507 vs FL$=$0.379), while FL edges ahead on overall MCC.

\begin{table*}[t]
\centering
\caption{Best performing configurations on the Flavor dataset. Prec and Rec are macro-averaged. AUC is one-vs-rest macro ROC AUC. $^{*}$, $^+$ denotes the top and second-best performing method by average metric.}
\label{tab:flavor}
\resizebox{\textwidth}{!}{%
\begin{tabular}{llllrrrrrrr}
\toprule
Modality & CLM & Loss & ML & AUC & Acc & F1$_\text{mac}$ & F1$_\text{mic}$ & MCC & Prec & Rec \\
\midrule
E2E LoRA$^+$ & MolFormer-XL & WL & --- & \textbf{0.975} & 0.892 & \textbf{0.809} & 0.892 & 0.801 & 0.802 & 0.817 \\
E2E LoRA & ChemBERTa-77M MTR & WL & --- & 0.964 & 0.890 & 0.780 & 0.890 & 0.802 & 0.780 & 0.786 \\
E2E LoRA & Uni-Mol & FL & --- & 0.969 & 0.883 & 0.771 & 0.883 & 0.790 & 0.767 & 0.781 \\
E2E LoRA & Uni-Mol & WL & --- & 0.973 & 0.885 & 0.783 & 0.885 & 0.795 & 0.747 & \textbf{0.848} \\
\midrule
FT Embed & MolFormer-XL & WL & LightGBM & 0.970 & 0.882 & 0.783 & 0.882 & 0.787 & 0.800 & 0.778 \\
FT Embed & Uni-Mol & FL & XGBoost & 0.950 & 0.866 & 0.731 & 0.866 & 0.765 & 0.756 & 0.735 \\
PC + FT Embed & MolFormer-XL & WL & LightGBM & 0.967 & 0.883 & 0.798 & 0.883 & 0.789 & 0.852 & 0.779 \\
\midrule
RAFE & ChemBERTa-77M MTR & FL & CatBoost & 0.969 & 0.886 & 0.741 & 0.886 & 0.798 & 0.731 & 0.756 \\
RAFE & MolFormer-XL & FL & LightGBM & 0.953 & 0.889 & 0.788 & 0.889 & 0.793 & 0.818 & 0.768 \\
RAFE & MolFormer-XL & WL & LightGBM & 0.964 & 0.892 & 0.794 & 0.892 & 0.798 & 0.825 & 0.775 \\
RAFE$^*$ & MolFormer-XL & WL & CatBoost & 0.963 & \textbf{0.893} & 0.774 & \textbf{0.893} & \textbf{0.804} & \textbf{0.868} & 0.749 \\
RAFE & Uni-Mol & FL & CatBoost & 0.931 & 0.866 & 0.714 & 0.866 & 0.764 & 0.704 & 0.731 \\
\bottomrule
\end{tabular}
}
\end{table*}

\subsection{BACE Dataset Results}
Table~\ref{tab:bace} compares key configurations on the BACE inhibitor prediction task. BACE is the most challenging dataset (MCC $< 0.65$ across most methods), consistent with the difficulty of generalizing to novel chemical scaffolds in a binding-affinity prediction task.

\textbf{Uni-Mol with weighted loss achieves the best overall performance} (MCC$=$0.623 $95\%$ CI $[0.546, 0.700]$, SE$=$0.039; AUC$=$0.891 $95\%$ CI $[0.841, 0.941]$; Acc$=$0.822), representing an 8\% absolute MCC gain over the next best E2E LoRA MolFormer FL (MCC$=$0.577, AUC$=$0.872); the CIs overlap (MolFormer FL 95\,\% CI $\approx [0.499, 0.656]$), so this gain is directionally consistent but does not reach statistical significance on $n{=}152$ test set. This improvement is consistent with the hypothesis that 3D geometry-aware representations provide superior inductive biases for binding affinity tasks: BACE-1 inhibitor activity is driven by ligand-receptor complementarity in three-dimensional space.

The RAFE framework with MolFormer-XL WL + LightGBM achieves MCC$=$0.561, intermediate between the finetuned-embedding-only result (MCC$=$0.557) and E2E LoRA (MCC$=$0.577). Adding PC features to finetuned embeddings provides minimal improvement over RAFE alone, consistent with the finding that LoRA-adapted representations already capture task-relevant molecular properties. Uni-Mol FT Embed (FL + LightGBM) achieves MCC$=$0.493, AUC$=$0.860, which is substantially below Uni-Mol E2E LoRA (MCC$=$0.623), confirming that end-to-end 3D gradient flow is
essential for binding affinity prediction and cannot be fully recovered at the embedding stage. Uni-Mol RAFE (FL + XGBoost, MCC$=$0.478, AUC$=$0.840) modestly improves over Uni-Mol FT Embed but remains well below Uni-Mol E2E LoRA, reaffirming that geometry-critical tasks benefit most from fully differentiable 3D encoding.

Figure~\ref{fig:bace} shows ROC and precision-recall curves for the best configuration per model family. MolFormer-XL base embeddings with CatBoost achieve AUC$=$0.885 (highest among SMILES-based approaches), while Uni-Mol WL achieves the highest AUPR$=$0.924 and overall MCC.

\begin{figure*}[!ht]
  \centering
  \includegraphics[width=0.49\textwidth]{Figures/bace_roc.pdf}
  \includegraphics[width=0.49\textwidth]{Figures/bace_pr.pdf}
  \caption{ROC and precision-recall curves for the BACE inhibitor prediction task. (a) Curves shown for the best configuration per model family (MolFormer E2E LoRA FL, Uni-Mol E2E LoRA WL, MolFormer RAFE). Uni-Mol WL achieves the highest AUPR~$= 0.924$ and overall MCC~$= 0.623$, while MolFormer base embeddings with CatBoost reach
  AUC~$= 0.885$ among SMILES-based approaches.}
  \label{fig:bace}
\end{figure*}

\begin{table*}[ht]
\centering
\caption{Best configurations on the BACE dataset. $^*$ overall best by MCC. Prec and Rec are macro-averaged.}
\label{tab:bace}
\resizebox{\textwidth}{!}{%
\begin{tabular}{llllrrrrrrr}
\toprule
Modality & CLM & Loss & ML & AUC & Acc & F1$_\text{mac}$ & F1$_\text{mic}$ & MCC & Prec & Rec \\
\midrule
E2E LoRA & MolFormer-XL & FL & --- & 0.872 & 0.789 & 0.785 & 0.789 & 0.577 & 0.783 & 0.794 \\
E2E LoRA & MolFormer-XL & WL & --- & 0.871 & 0.724 & 0.722 & 0.724 & 0.476 & 0.733 & 0.743 \\
E2E LoRA & Uni-Mol & FL & --- & 0.876 & 0.743 & 0.743 & 0.743 & 0.520 & 0.755 & 0.765 \\
E2E LoRA$^*$ & Uni-Mol & WL & --- & \textbf{0.891} & \textbf{0.822} & \textbf{0.808} & \textbf{0.822} & \textbf{0.623} & \textbf{0.823} & \textbf{0.801} \\
\midrule
FT Embed & MolFormer-XL & FL & CatBoost & 0.869 & 0.763 & 0.762 & 0.763 & 0.557 & 0.773 & 0.784 \\
FT Embed & Uni-Mol & FL & LightGBM & 0.860 & 0.730 & 0.729 & 0.730 & 0.493 & 0.742 & 0.751 \\
PC + FT Embed & MolFormer-XL & FL & XGBoost & 0.870 & 0.737 & 0.736 & 0.737 & 0.517 & 0.755 & 0.762 \\
\midrule
RAFE & MolFormer-XL & FL & CatBoost & 0.863 & 0.763 & 0.762 & 0.763 & 0.557 & 0.773 & 0.784 \\
RAFE & MolFormer-XL & WL & LightGBM & 0.858 & 0.770 & 0.768 & 0.770 & 0.561 & 0.774 & 0.787 \\
RAFE & Uni-Mol & FL & XGBoost & 0.840 & 0.704 & 0.704 & 0.704 & 0.478 & 0.739 & 0.738 \\
\bottomrule
\end{tabular}
}
\end{table*}

\subsection{ClinTox Dataset Results}
Table~\ref{tab:clintox} presents results for the FDA\_APPROVED label. \textbf{RAFE with MolFormer-XL FL + CatBoost achieves the overall best performance} (MCC$=$0.913 $95\%$ CI $[0.867, 0.959]$, SE$=$0.024; AUC$=$ 0.986 $95\%$ CI $[0.967, 1.000]$; Acc$=$0.986, F1$_\text{mac}$=0.955), demonstrating that ZINC-250k neighborhood context provides a strong complementary signal for FDA approval classification. E2E LoRA (MolFormer FL) achieves MCC$=$0.786, while ChemBERTa MTR WL LoRA reaches MCC$=$0.750. Uni-Mol E2E LoRA (FL, MCC$=$ 0.282) performs substantially below SMILES-based models, indicating that 3D geometric features cannot encapsulate meaningful representation for FDA approval classification.

Among non-RAFE approaches, FT~Embed with ChemBERTa-77M MLM FL + CatBoost (MCC$=$0.850, second-best overall) shows that ChemBERTa's masked-language pre-training captures FDA-approval-relevant chemical patterns that CatBoost can exploit effectively from frozen embeddings. Adding PC features to FT embeddings (PC + FT Embed, ChemBERTa MLM WL + LightGBM: MCC$=$0.786) reaches E2E LoRA parity without retrieval. RAFE (ChemBERTa MLM WL + LightGBM: MCC$=$0.869 $95\%$ CI $[0.814, 0.924]$; MolFormer FL + CatBoost: MCC$=$0.913 $95\%$ CI $[0.867, 0.959]$) surpasses all non-retrieval baselines by adding ZINC-250k neighborhood context. Uni-Mol RAFE (WL + XGBoost: MCC$=$0.633) substantially
outperforms Uni-Mol E2E LoRA (MCC$=$0.282), demonstrating that ZINC-250k retrieval compensates for 3D geometry's limited transferability to toxicity prediction.

\begin{table*}[!ht]
\centering
\caption{Best configurations per modality on the ClinTox FDA\_APPROVED label.
$^*$ overall best by MCC; $^+$ second-best. Prec and Rec are macro-averaged. Rec estimated from macro F1 and Prec where not directly recorded ($^\dagger$).}
\label{tab:clintox}
\resizebox{\textwidth}{!}{%
\begin{tabular}{llll rrrrrrrr}
\toprule
Modality & CLM & Loss & ML & AUC & Acc & F1$_\text{mac}$ & F1$_\text{mic}$ & MCC & Prec & Rec \\
\midrule
E2E LoRA & ChemBERTa-77M MLM & FL & --- & 0.987 & 0.958 & 0.839 & 0.958 & 0.717 & 0.769 & 0.978 \\
E2E LoRA & ChemBERTa-77M MTR & WL & --- & 0.987 &  0.965 & 0.859 & 0.965 & 0.750 & 0.792 & 0.982 \\
E2E LoRA & MolFormer-XL & FL & --- & \textbf{0.993} &  0.972 & 0.881 & 0.972 & 0.786 & 0.818 & 0.985 \\
E2E LoRA & Uni-Mol & FL & --- & 0.904 & 0.867 & 0.612 & 0.867 & 0.282 & 0.588 & 0.727 \\
\midrule
FT Embed$^+$ & ChemBERTa-77M MLM & FL & CatBoost & 0.987 & 0.986 & 0.925 & 0.986 & 0.850 & 0.925 & 0.925 \\
FT Embed & MolFormer-XL & FL & LightGBM & 0.987 & 0.979 & 0.906 & 0.979 & 0.827 & 0.850 & \textbf{0.989} \\
FT Embed & Uni-Mol & FL & CatBoost & 0.923 & 0.853 & 0.658 & 0.853 & 0.460 & 0.625 & 0.923 \\
\midrule
PC + FT Embed & ChemBERTa-77M MLM & WL & LightGBM & 0.989 & 0.972 & 0.881 & 0.972 & 0.786 & 0.818 & 0.985 \\
PC + FT Embed & MolFormer-XL & FL & CatBoost & 0.991 & 0.958 & 0.839 & 0.958 & 0.717 & 0.769 & 0.978 \\
\midrule
RAFE & ChemBERTa-77M MLM & WL & LightGBM & 0.989 & 0.979 & 0.934 & 0.979 & 0.869 & 0.951 & 0.917$^\dagger$ \\
RAFE$^*$ & MolFormer-XL & FL & CatBoost & 0.986 & \textbf{0.986} & \textbf{0.955} & \textbf{0.986} & \textbf{0.913} & \textbf{0.992} & 0.921$^\dagger$ \\
RAFE & Uni-Mol & WL & XGBoost & 0.983 & 0.951 & 0.803 & 0.951 & 0.633 & 0.746 & 0.907 \\
\bottomrule
\end{tabular}
}
\end{table*}

\subsection{Parameter Efficiency}
Table~\ref{tab:params} reports LoRA hyperparameters and trainable parameter counts for all model families across the four tasks. ChemBERTa variants (3.4M base parameters) achieve reductions of $76$-$94\%$ depending on rank; MolFormer-XL ($\approx 46$M parameters) achieves $84$-$97\%$ reduction. For Uni-Mol (rank $r \in \{4,8,16,32\}$), LoRA targets only the \texttt{in\_proj} weight in $15$ attention layers, yielding a particularly efficient adaptation with $<5\%$ trainable parameters for the lowest-rank configurations.

%Full parameter counts are tabulated in Supplementary Table~S6.
\begin{table*}[!ht]
\centering
\caption{Selected LoRA configurations illustrating the range of parameter efficiency. Trainable~\% = 100 $\times$ trainable / total parameters.}
\label{tab:params}
\resizebox{\columnwidth}{!}{%
\begin{tabular}{llllrrrr}
\toprule
Dataset & Model & $r$ & $\alpha$ & Total & Trainable & Trainable \% & \% Reduction \\
\midrule
BACE   & ChemBERTa-MLM & 4  & 32  & 3.43M & 205K  & 5.66  & 94.34 \\
BACE   & ChemBERTa-MTR & 32 & 128 & 3.43M & 606K  & 15.03 & 84.97 \\
BACE   & MolFormer-XL  & 64 & 128 & 45.6M & 8.26M & 15.35 & 84.65 \\
BBBP   & ChemBERTa-MLM & 64 & 32  & 3.43M & 1.06M & 23.69 & 76.31 \\
BBBP   & MolFormer-XL  & 16 & 16  & 45.6M & 2.95M & 6.09  & 93.91 \\
ClinTox& MolFormer-XL  & 4  & 32  & 45.6M & 1.63M & 3.44  & 96.56 \\
Flavor & MolFormer-XL  & 4  & 64  & 45.6M & 1.63M & 3.44  & 96.56 \\
\bottomrule
\end{tabular}
}
\end{table*}

\subsection{Comparison with External Baselines}\label{sec:external_comparison}
Table~\ref{tab:external_main} compares the best PEARL configurations against published
graph-based and language-model baselines on the same MoleculeNet scaffold splits.
AUC is the primary comparison metric as it is most consistently reported across sources;
full method descriptions and the ClinTox split-difference caveat are in Supplementary Section~S5.

\begin{table*}[!ht]
\centering
\caption{AUC comparison with published external baselines on MoleculeNet scaffold splits.
\textsuperscript{a}~Mean~$\pm$~std over multiple seeds (cited source).
\textsuperscript{b}~PEARL point estimate on the same splits.
``-''~= not reported by source.}
\label{tab:external_main}
\resizebox{\columnwidth}{!}{%
\begin{tabular}{llccc}
\toprule
\textbf{Method} & \textbf{Ref.} & \textbf{BACE AUC} & \textbf{BBBP AUC} & \textbf{Flavor AUC} \\
\midrule
\multicolumn{5}{l}{\textit{Graph neural network baselines}} \\[2pt]
Chemprop (D-MPNN)      & \cite{yang2019chemprop}  & $0.861 \pm 0.034$\textsuperscript{a} & $0.919 \pm 0.026$\textsuperscript{a} & --- \\
GROVER-Large           & \cite{rong2020grover}    & $0.894 \pm 0.028$\textsuperscript{a} & $0.940 \pm 0.016$\textsuperscript{a} & --- \\
Mole-BERT              & \cite{xia2023molecbert}  & $0.856 \pm 0.029$\textsuperscript{a} & $0.940 \pm 0.011$\textsuperscript{a} & --- \\
\midrule
\multicolumn{5}{l}{\textit{Chemical language models (full finetuning)}} \\[2pt]
Uni-Mol (full FT)      & \cite{zhou2023unimol}    & $0.857 \pm 0.002$\textsuperscript{a} & $0.932 \pm 0.003$\textsuperscript{a} & --- \\
MolFormer-XL (full FT) & \cite{ross2022molformer} & ---                                  & $0.932$\textsuperscript{a}           & --- \\
Flavor model             & \cite{zimmermann2024fart}& ---                                  & ---                                  & $0.964$ \\
\midrule
\multicolumn{5}{l}{\textit{PEARL (this work, LoRA finetuning)}} \\[2pt]
PEARL Uni-Mol WL (E2E LoRA)   & this work & $\mathbf{0.891}$\textsuperscript{b} & $0.943$\textsuperscript{b}          & $0.973$\textsuperscript{b} \\
PEARL MolFormer WL (E2E LoRA) & this work & $0.871$\textsuperscript{b}          & $\mathbf{0.948}$\textsuperscript{b} & $\mathbf{0.975}$\textsuperscript{b} \\
PEARL RAFE best               & this work & $0.875$\textsuperscript{b}          & $0.940$\textsuperscript{b}          & $0.970$\textsuperscript{b} \\
\bottomrule
\end{tabular}
}
\end{table*}

On \textbf{BACE}, PEARL Uni-Mol WL (AUC~$= 0.891$) is competitive with GROVER-Large ($0.894 \pm 0.028$) within the reported standard deviation, and surpasses the published Uni-Mol full-finetuning result ($0.857$)~\cite{zhou2023unimol}, demonstrating that LoRA regularisation is beneficial on this 1{,}513-molecule dataset. On \textbf{BBBP}, PEARL MolFormer WL E2E LoRA (AUC~$= 0.948$) exceeds all GNN baselines including GROVER-Large ($0.940$) and Mole-BERT ($0.940$), and surpasses MolFormer full finetuning ($0.932$)~\cite{ross2022molformer} while using only $6\%$ of its trainable
parameters. On the \textbf{Flavor} multi-class task, PEARL MolFormer WL (AUC~$= 0.975$) exceeds the task-specific FART language model ($0.964$)~\cite{zimmermann2024fart}. For \textbf{ClinTox}, PEARL E2E LoRA (AUC~$= 0.993$) and RAFE MolFormer FL
(AUC~$= 0.986$, MCC~$= 0.913$) substantially exceed published GNN and CLM baselines
under our evaluation protocol; a split-difference caveat applies (see Supplementary~S5).
%%%%%%%%%%%%%%%%%%%%%%%%%%%%%%%%%%%%%%%%%%%%%%%%%%%%%%%%%%%%%%%%

\section{Discussion \& Conclusion}
PEARL extends parameter-efficient molecular property prediction in two directions that address limitations of SMILES-only, single-modality CLM finetuning. 

\textbf{3D geometry matters for binding-affinity tasks:}
Uni-Mol's SE(3)-equivariant architecture provides the most accurate BACE inhibitor classification (MCC$=$0.623 $95\%$ CI $[0.546, 0.700]$, vs.\ MCC$=$0.577 $95\%$ CI $[0.499, 0.656]$ for the best SMILES-based approach). BACE-1 inhibitor activity is governed by three-dimensional ligand-receptor complementarity, and the inductive biases of a geometry-aware transformer are therefore well-matched to this task. The hybrid 2560-dimensional representation (512-dim 3D CLS token plus 2048-dim ECFP4) balances 3D structural awareness with explicit 2D topological information. Conversely, Uni-Mol does not improve on SMILES-based models for ClinTox, where the primary challenge is severe class imbalance rather than geometric discrimination.

\textbf{Retrieval-augmented chemical neighborhood context provides orthogonal information:}
The RAFE framework demonstrates that ZINC-250k neighborhood features such as logP/QED/SAS distributions, neighborhood density and isolation, embedding centroid geometry, and Tanimoto structural similarity, provides information complementary to learned model representations. On BBBP, RAFE achieves MCC$=$ 0.798 ($95\%$ CI $[0.743, 0.853]$), essentially matching E2E LoRA (MCC$=$0.797) with a fully interpretable tree-based pipeline; the overlapping CIs confirm these two approaches are statistically indistinguishable on this dataset. On ClinTox FDA\_APPROVED, RAFE (ChemBERTa MLM WL + LightGBM: MCC$=$0.869 $95\%$ CI $[0.814, 0.924]$; MolFormer FL + CatBoost: MCC$=$ 0.913 $95\%$ CI $[0.867, 0.959]$) substantially improves over both E2E LoRA (MolFormer FL: MCC$=$0.786) and PC + FT Embed (MCC$=$0.786) by supplying ZINC-250k neighborhood context for approval prediction. The non-overlapping CIs between RAFE MolFormer FL + CatBoost ($95\%$ CI $[0.867, 0.959]$) and E2E LoRA ($95\%$ CI $[0.719, 0.853]$) confirm this improvement is statistically significant. RAFE MolFormer FL + CatBoost (MCC$=$0.913) also surpasses the best non-retrieval configuration, FT Embed with ChemBERTa MLM FL + CatBoost (MCC$=$0.850), indicating that ZINC-250k neighborhood context provides drug-likeness and physicochemical signals complementary to the sequence-level patterns encoded in frozen embeddings. The superiority of CatBoost over LightGBM and XGBoost for this RAFE configuration is consistent with CatBoost's robustness to high-cardinality, mixed-type feature vectors. Uni-Mol RAFE (WL + XGBoost: MCC$=$0.633) nonetheless substantially outperforms Uni-Mol E2E LoRA (MCC$=$0.282) for ClinTox, confirming that retrieval augmentation compensates effectively when end-to-end 3D learning fails under severe class imbalance.

\textbf{Task-dependent optimal strategies persist:} Several consistent patterns emerge across the expanded study:
(i)~MolFormer consistently outperforms ChemBERTa across all experimental conditions and tasks; 
(ii)~weighted loss excels on severely imbalanced datasets (BBBP, Flavor), while focal
loss is optimal for challenging tasks (ClinTox); 
(iii)~adding PC features to finetuned embeddings provides minimal improvement over finetuned embeddings alone, because LoRA adaptation effectively internalizes task-relevant molecular properties;
(iv)~for ClinTox, base embeddings combined with PC features remain a strong baseline, and RAFE is needed to outperform them with finetuned representations to achieve best Recall.

\textbf{Parameter efficiency:}
LoRA achieves $75$-$97\%$ parameter reduction across all three model families without sacrificing performance. The smallest effective configurations (MolFormer-XL at $r = 4$, $\alpha = 64$ for Flavor and ClinTox, $3.44\%$ trainable) are particularly noteworthy: they demonstrate that the intrinsic task-relevant subspace of a $46$M-parameter molecular transformer can be captured with fewer than 1.7M additional parameters.

\textbf{Future directions:}
Some avenues still remain open-ended. First, extending PEARL to regression tasks (binding affinity, solubility) with calibrated uncertainty estimates would broaden its applicability. Secondly, emerging multi-modal CLMs such as nach0 \cite{livne2024nach0} and Omni-Mol \cite{hu2025omnimol}, which integrate SMILES, graphs, and textual descriptions, represent natural extensions of the LoRA + RAFE paradigm. 

In conclusion, PEARL establishes a practical framework for molecular property prediction that combines parameter-efficient 3D and SMILES-based finetuning with retrieval-augmented chemical neighborhood context. By achieving state-of-the-art performance on ClinTox toxicity prediction (RAFE MCC$=$0.913 $95\%$ CI $[0.867, 0.959]$) and Flavor multi-class prediction (RAFE MCC$=$0.804 $95\%$ CI $[0.784, 0.824]$), and competitive performance on BACE inhibitor classification (Uni-Mol MCC$=$0.623) and BBBP (RAFE MCC$=$0.798 $95\%$ CI $[0.743, 0.853]$), all while reducing trainable parameters by $75$-$97\%$, PEARL enables resource-constrained laboratories to deploy high-performance molecular AI tools.
%%%%%%%%%%%%%%%%%%%%%%%%%%%%%%%%%%%%%%%%%%%

\section*{Key Points}
\begin{itemize}
  \item PEARL applies LoRA to three CLM families - ChemBERTa, MolFormer, and Uni-Mol, achieving $75$-$97\%$ parameter reduction while matching or exceeding full finetuning.
  \item Uni-Mol's SE(3)-equivariant 3D representations achieve the best BACE inhibitor classification (MCC$=$0.623 $95\%$ CI $[0.546, 0.700]$); the 8\% absolute MCC gain over the best SMILES-based method (MCC$=$ 0.577).
  \item MolFormer consistently outperforms ChemBERTa; weighted loss excels on imbalanced datasets while focal loss is optimal for challenging tasks like ClinTox.
  \item Adding PC features to finetuned embeddings provides minimal benefit, as LoRA-adapted models already internalize task-relevant molecular properties.
\end{itemize}

\section*{Declarations}
\begin{itemize}
\item Funding Sources

This research did not receive any specific grant from funding agencies in the public, commercial, or not-for-profit sectors.

\item Competing interests

The authors declare no competing financial interests. R.M.\ is Scientific Advisor at Eternal, U.A.E.

\item Data availability 

All code and data associated with PEARL is available at \url{https://github.com/raghvendra5688/PEARL}.

\item Acknowledgments

We thank the Qatar National Library (QNL) for their support in providing resources and funding the APC.

\item Author contributions

K.A. contributed to Methods, Experiments, Validation, Writing – original draft. S.G. and H.B. participated in Project Supervision, Writing - review \& editing. R.M. Conceptualized, Supervised the project and contributed in Methods, Experiments, Validation, Writing - review \& editing.
\end{itemize}

\bmhead{Supplementary information}

The supplementary material accompanying this article contains additional experimental details, complete ablation studies, feature descriptions, and implementation details supporting the results presented in the main manuscript.

\begin{appendices}

\section{Supplementary Material}\label{sec:supplementary}

\noindent\textbf{Supplementary Table S1.} RDKit-derived molecular descriptors, NetworkX graph
features, Morgan fingerprints, and MACCS keys comprising the 473-dimensional engineered
feature vector.

\noindent\textbf{Supplementary Table S2.} Full BBBP ablation: all modality $\times$ CLM $\times$
loss $\times$ ML combinations.

\noindent\textbf{Supplementary Table S3.} Full Flavor ablation including per-class AUPR.

\noindent\textbf{Supplementary Table S4.} Full ClinTox ablation for FDA\_APPROVED.

\noindent\textbf{Supplementary Table S5.} Full BACE ablation.

\noindent\textbf{Supplementary Table S6.} Complete LoRA parameter efficiency analysis for all
model families, datasets, and loss functions, including Uni-Mol configurations.

\section{Additional Methodological Details}\label{sec:supp_methods}

\noindent\textbf{Supplementary Section S1.} Zero-shot vs.\ full finetuning: background and
comparison.

\noindent\textbf{Supplementary Section S2.} Detailed LoRA mathematical formulation and
implementation for Uni-Mol (\texttt{in\_proj} targeting, gradient flow preservation,
conformer caching strategy).

\noindent\textbf{Supplementary Section S3.} PEARL custom trainer: loss function
implementations, class-weight computation, and evaluation metric computation.

\noindent\textbf{Supplementary Section S4.} RAFE pipeline details: ZINC-250k cleaning
procedure, FAISS index construction (FlatIP vs.\ HNSW trade-offs), PCA fitting strategy,
and per-group feature breakdown for the 69-dimensional RAFE vector.

\end{appendices}
%%%%%%%%%%%%%%%%%%%%%%%%%%%%%%%%%%%%%%%%%%%%%%%%%%%%%%

\bibliography{sn-bibliography}
% common bib file
%% if required, the content of .bbl file can be included here once bbl is generated
%%\input sn-article.bbl

\end{document}
